\documentclass[a4paper,12pt]{ctexart}
\usepackage{../../teaching-style}

\begin{document}
\pagestyle{empty}

\maintitle{微积分}
\begin{center}
  \large 高学段 · 导数与积分
\end{center}
\vspace{0.3cm}

\chaptertitle{第一章\quad 极限}

\sectiontitle{1.1\quad 数列的极限}

\knowbox{
数列$\{a_n\}$的极限：当$n$无限增大时，$a_n$无限趋近于常数$A$，记作$\lim\limits_{n\to\infty}a_n=A$。\\
常见极限：$\lim\limits_{n\to\infty}\dfrac1n=0$，$\lim\limits_{n\to\infty}q^n=0$（$|q|<1$）。\\
运算法则：$\lim(a_n\pm b_n)=\lim a_n\pm\lim b_n$，$\lim(a_nb_n)=\lim a_n\cdot\lim b_n$，\\
$\lim\dfrac{a_n}{b_n}=\dfrac{\lim a_n}{\lim b_n}$（分母$\neq0$）。}

\sectiontitle{1.2\quad 函数的极限与连续}

\knowbox{
$\lim\limits_{x\to x_0}f(x)=A$：当$x$无限趋近$x_0$时，$f(x)$无限趋近$A$。\\
$\lim\limits_{x\to\infty}f(x)=A$：当$|x|$无限增大时，$f(x)$趋近$A$。\\
\textbf{连续}：$\lim\limits_{x\to x_0}f(x)=f(x_0)$。初等函数在其定义域内连续。\\
\textbf{两个重要极限}：$\lim\limits_{x\to0}\dfrac{\sin x}{x}=1$，$\lim\limits_{x\to\infty}\left(1+\dfrac1x\right)^x=e$。}

\sectiontitle{习题1}

\begin{exercise}
\item 求极限：$\lim\limits_{n\to\infty}\dfrac{3n+1}{2n-1}$；$\lim\limits_{n\to\infty}\dfrac{2n^2+3}{n^2+1}$。
\item 求极限：$\lim\limits_{x\to2}(3x^2-2x+1)$；$\lim\limits_{x\to1}\dfrac{x^2-1}{x-1}$。
\item 求极限：$\lim\limits_{x\to0}\dfrac{\sin3x}{x}$；$\lim\limits_{x\to0}(1+2x)^{\frac1x}$。
\end{exercise}

\newpage

\chaptertitle{第二章\quad 导数}

\sectiontitle{2.1\quad 导数的概念}

\knowbox{
$f'(x_0)=\lim\limits_{\Delta x\to0}\dfrac{f(x_0+\Delta x)-f(x_0)}{\Delta x}$。\\
几何意义：$f'(x_0)$是曲线$y=f(x)$在点$(x_0,f(x_0))$处切线的斜率。\\
物理意义：$s'(t)$是瞬时速度，$v'(t)$是加速度。}

\sectiontitle{2.2\quad 基本导数公式与运算法则}

\knowbox{
$(c)'=0$，$(x^n)'=nx^{n-1}$，$(e^x)'=e^x$，$(\ln x)'=\dfrac1x$，\\
$(\sin x)'=\cos x$，$(\cos x)'=-\sin x$。\\
四则运算法则：$(u\pm v)'=u'\pm v'$，$(uv)'=u'v+uv'$，$\left(\dfrac{u}{v}\right)'=\dfrac{u'v-uv'}{v^2}$。\\
链式法则：$[f(g(x))]'=f'(g(x))\cdot g'(x)$。\\
隐函数求导：将$y$视为$x$的函数，两边同时对$x$求导。}

\sectiontitle{2.3\quad 导数的应用}

\knowbox{
\textbf{单调性}：$f'(x)>0\Rightarrow$增，$f'(x)<0\Rightarrow$减。\\
\textbf{极值}：$f'(x_0)=0$，左正右负为极大值，左负右正为极小值。\\
\textbf{最值}：比较极值点和端点的函数值。\\
\textbf{切线}：$y-f(x_0)=f'(x_0)(x-x_0)$。\\
\textbf{微分中值定理}：若$f$在$[a,b]$连续、$(a,b)$可导，则存在$\xi\in(a,b)$使$f(b)-f(a)=f'(\xi)(b-a)$。}

\sectiontitle{习题2}

\begin{exercise}
\item 求导数：$y=x^3-3x^2+2x-1$；$y=(2x+1)^3$；$y=\ln(x^2+1)$。
\item 求$y=x^3-3x+1$的单调区间和极值。
\item 求$y=x^3-3x^2+2$在$[-1,3]$上的最值。
\item 求$y=x^2$在$(1,1)$处的切线方程。
\item 证明：$x>0$时，$x-\dfrac{x^3}{6}<\sin x<x$（用导数与凹凸性）。
\end{exercise}

\newpage

\chaptertitle{第三章\quad 微分与泰勒展开}

\sectiontitle{3.1\quad 微分}

\knowbox{
$dy=f'(x)dx$。\\
近似公式：$f(x+\Delta x)\approx f(x)+f'(x)\Delta x$。}

\sectiontitle{3.2\quad 泰勒展开}

\knowbox{
$f(x)=f(x_0)+f'(x_0)(x-x_0)+\dfrac{f''(x_0)}{2!}(x-x_0)^2+\cdots+\dfrac{f^{(n)}(x_0)}{n!}(x-x_0)^n+R_n$。\\
\textbf{常用展开}（在$x=0$处，麦克劳林公式）：\\
$e^x=1+x+\dfrac{x^2}{2!}+\dfrac{x^3}{3!}+\cdots$\\
$\sin x=x-\dfrac{x^3}{3!}+\dfrac{x^5}{5!}-\cdots$\\
$\cos x=1-\dfrac{x^2}{2!}+\dfrac{x^4}{4!}-\cdots$\\
$\ln(1+x)=x-\dfrac{x^2}{2}+\dfrac{x^3}{3}-\cdots$\\
$(1+x)^\alpha=1+\alpha x+\dfrac{\alpha(\alpha-1)}{2!}x^2+\cdots$}

\sectiontitle{习题3}

\begin{exercise}
\item 用微分求$\sqrt{4.01}$的近似值。
\item 写出$e^x$的$4$阶麦克劳林展开。
\item 写出$\sin x$的$5$阶麦克劳林展开，并用它估算$\sin0.1$。
\end{exercise}

\newpage

\chaptertitle{第四章\quad 不定积分}

\sectiontitle{4.1\quad 原函数与不定积分}

\knowbox{
若$F'(x)=f(x)$，则$\int f(x)\,dx=F(x)+C$。\\
基本积分公式：$\int x^n\,dx=\dfrac{x^{n+1}}{n+1}+C$（$n\neq-1$），\\
$\int\dfrac1x\,dx=\ln|x|+C$，$\int e^x\,dx=e^x+C$，\\
$\int\sin x\,dx=-\cos x+C$，$\int\cos x\,dx=\sin x+C$。\\
线性性质：$\int[af(x)+bg(x)]\,dx=a\int f(x)\,dx+b\int g(x)\,dx$。}

\sectiontitle{4.2\quad 换元积分法}

\knowbox{
\textbf{第一类换元}（凑微分）：$\int f(g(x))g'(x)\,dx=\int f(u)\,du$（$u=g(x)$）。\\
\textbf{第二类换元}：令$x=\varphi(t)$，$\int f(x)\,dx=\int f(\varphi(t))\varphi'(t)\,dt$。\\
常见凑微分：$\int f(ax+b)\,dx=\dfrac1a\int f(ax+b)\,d(ax+b)$。}

\sectiontitle{4.3\quad 分部积分法}

\knowbox{$\int u\,dv=uv-\int v\,du$。\\
选择原则：幂对优先（指数、三角函数容易积分时放在$dv$中）。}

\sectiontitle{习题4}

\begin{exercise}
\item 求不定积分：$\int(x^3-2x+1)\,dx$；$\int\dfrac1{x^2}\,dx$；$\int\sqrt{x}\,dx$。
\item 用凑微分法：$\int(2x+1)^5\,dx$；$\int\dfrac{x}{x^2+1}\,dx$；$\int\sin3x\,dx$。
\item 用分部积分法：$\int xe^x\,dx$；$\int x\sin x\,dx$；$\int\ln x\,dx$。
\end{exercise}

\newpage

\chaptertitle{第五章\quad 定积分}

\sectiontitle{5.1\quad 定积分的概念}

\knowbox{
$\int_a^b f(x)\,dx=\lim\limits_{\lambda\to0}\sum_{i=1}^n f(\xi_i)\Delta x_i$（黎曼和）。\\
几何意义：$x$轴上方曲边梯形的面积减去下方的面积。\\
\textbf{微积分基本定理（牛顿-莱布尼茨公式）}：\\
若$F'(x)=f(x)$，则$\int_a^b f(x)\,dx=F(b)-F(a)$。}

\sectiontitle{5.2\quad 定积分的计算与应用}

\knowbox{
\textbf{计算}：\
换元法：$\int_a^b f(\varphi(x))\varphi'(x)\,dx=\int_{\varphi(a)}^{\varphi(b)} f(u)\,du$。\\
分部法：$\int_a^b u\,dv=[uv]_a^b-\int_a^b v\,du$。\\
\textbf{应用}：\\
面积：$S=\int_a^b|f(x)-g(x)|\,dx$。\\
旋转体体积：$V=\pi\int_a^b[f(x)]^2\,dx$（绕$x$轴）。\\
弧长：$L=\int_a^b\sqrt{1+[f'(x)]^2}\,dx$。\\
物理应用：变力做功$W=\int_a^b F(x)\,dx$。}

\sectiontitle{习题5}

\begin{exercise}
\item 计算：$\int_0^1 (x^2+1)\,dx$；$\int_0^\pi\sin x\,dx$；$\int_1^2\dfrac1x\,dx$。
\item 求$y=x^2$与$y=x$所围成图形的面积。
\item 求$y=\sqrt{x}$在$[0,4]$上绕$x$轴旋转的体积。
\item 求$y=x^{\frac23}$在$[1,8]$上的弧长。
\end{exercise}

\newpage

\chaptertitle{第六章\quad 常微分方程初步}

\sectiontitle{6.1\quad 可分离变量方程}

\knowbox{
形如$\dfrac{dy}{dx}=f(x)g(y)$，分离变量：$\dfrac{dy}{g(y)}=f(x)dx$，两边积分。}

\sectiontitle{6.2\quad 一阶线性方程}

\knowbox{
形如$y'+P(x)y=Q(x)$，通解公式：\\
$y=e^{-\int P\,dx}\left(\int Q e^{\int P\,dx}\,dx+C\right)$。}

\sectiontitle{习题6}

\begin{exercise}
\item 解微分方程：$\dfrac{dy}{dx}=2xy$。
\item 解微分方程：$y'=y$，且$y(0)=1$。
\item 解一阶线性方程：$y'+2xy=x$。
\end{exercise}

\newpage
\sectiontitle{参考答案}

\subsection*{习题1}
\begin{enumerate}
  \item $\lim\limits_{n\to\infty}\dfrac{3n+1}{2n-1}=\dfrac32$；$\lim\limits_{n\to\infty}\dfrac{2n^2+3}{n^2+1}=2$
  \item $\lim\limits_{x\to2}(3x^2-2x+1)=9$；$\lim\limits_{x\to1}\dfrac{x^2-1}{x-1}=2$
  \item $\lim\limits_{x\to0}\dfrac{\sin3x}{x}=3$；$\lim\limits_{x\to0}(1+2x)^{\frac1x}=e^2$
\end{enumerate}

\subsection*{习题2}
\begin{enumerate}
  \item $y'=3x^2-6x+2$；$y'=6(2x+1)^2$；$y'=\dfrac{2x}{x^2+1}$
  \item 增区间$(-\infty,-1)\cup(1,\infty)$，减区间$(-1,1)$；极大值$f(-1)=3$，极小值$f(1)=-1$
  \item 最大值$2$（$x=0$或$x=3$），最小值$-2$（$x=-1$或$x=2$）
  \item $y=2x-1$
  \item 设$f(x)=x-\sin x$，$f'(x)=1-\cos x\ge0$，$f(x)>0$得$\sin x<x$；设$g(x)=\sin x-(x-\frac{x^3}{6})$，利用$g'''(x)=1-\cos x\ge0$可证$g(x)>0$
\end{enumerate}

\subsection*{习题3}
\begin{enumerate}
  \item $\sqrt{4.01}\approx2.0025$
  \item $e^x=1+x+\dfrac{x^2}{2!}+\dfrac{x^3}{3!}+\dfrac{x^4}{4!}+O(x^5)$
  \item $\sin x=x-\dfrac{x^3}{3!}+\dfrac{x^5}{5!}+O(x^7)$；$\sin0.1\approx0.0998333$
\end{enumerate}

\subsection*{习题4}
\begin{enumerate}
  \item $\dfrac{x^4}{4}-x^2+x+C$；$-\dfrac1x+C$；$\dfrac23x^{\frac32}+C$
  \item $\dfrac{(2x+1)^6}{12}+C$；$\dfrac12\ln(x^2+1)+C$；$-\dfrac13\cos3x+C$
  \item $e^x(x-1)+C$；$-x\cos x+\sin x+C$；$x\ln x-x+C$
\end{enumerate}

\subsection*{习题5}
\begin{enumerate}
  \item $\dfrac43$；$2$；$\ln2$
  \item $\dfrac16$
  \item $8\pi$
  \item $\dfrac1{27}\bigl(80\sqrt{10}-13\sqrt{13}\bigr)$
\end{enumerate}

\subsection*{习题6}
\begin{enumerate}
  \item $y=Ce^{x^2}$
  \item $y=e^x$
  \item $y=\dfrac12+Ce^{-x^2}$
\end{enumerate}

\end{document}
